% Originally: content/week-09.tex, \section{Length, area, and volume} (Corsin Week 9)

\section{Length, area, and volume}
\label{sec:length_area_volume}

\subsection{The area of a parallelogram in \texorpdfstring{$\mathbb{R}^2$}{R2}}
\label{sec:area_parallelogram}

The relationship between the determinant and volume is a key concept in this part of the
lecture. We illustrate the special case of $\mathbb{R}^2$.

\begin{proposition}[Area of a parallelogram]
\label{prop:area_parallelogram_r2}
Let $x,y \in \mathbb{R}^2$. The area of the parallelogram $(0,x,x+y,y)$ is
\[A = \big|\det(x\mid y)\big|.\]
\end{proposition}

\begin{proof}
Assume w.l.o.g.\ that $x = x_1e_1$, a reduction justified below. Then, writing
$y = (y_1,y_2)\transp$,
\[A = |x_1y_2| = \left|\det\begin{pmatrix} x_1 & y_1 \\ 0 & y_2 \end{pmatrix}\right|,\]
directly from the base-times-height formula for a parallelogram with one side on the
$e_1$-axis.

% Sonnet 5 (Medium) --- FIG-W09-01
\begin{center}
\begin{tikzpicture}[>=Latex, scale=1.0]
  \draw[->] (-0.3,0) -- (3.6,0) node[right, font=\scriptsize] {$e_1$};
  \draw[->] (0,-0.3) -- (0,2.1) node[above, font=\scriptsize] {$e_2$};
  \coordinate (O) at (0,0);
  \coordinate (X) at (2.4,0);
  \coordinate (Y) at (0.7,1.7);
  \coordinate (XY) at (3.1,1.7);
  \draw[purple, dotted, thick] (O) -- (X) -- (XY) -- (Y) -- cycle;
  \draw[blue, thick, ->] (O) -- (X) node[below, font=\scriptsize] {$x=x_1e_1$};
  \draw[orange, thick, ->] (O) -- (Y) node[above left, font=\scriptsize] {$y$};
  \draw[green!50!black, thick, dotted] (Y) -- (0.7,0);
  \draw[green!50!black, ->] (0,-0.15) -- (0.7,-0.15) node[midway, below, font=\scriptsize] {$y_1$};
  \draw[green!50!black, ->] (0.85,0) -- (0.85,1.7) node[midway, right, font=\scriptsize] {$y_2$};
\end{tikzpicture}
\end{center}

For the reduction, let $R \in \operatorname{O}(2)$ be any rotation carrying $x,y$ to a pair
$Rx,Ry$ with $Rx$ along $e_1$; rotations preserve area, so
\[A = \big|\det(Rx\mid Ry)\big| = \big|\det(R)\det(x\mid y)\big| = \big|\det(x\mid y)\big|,\]
using $|\det R| = 1$ and multiplicativity of $\det$. So, the \qt{w.l.o.g.} above was
justified: the formula $A=|\det(x\mid y)|$ holds for the rotated pair, and equals
$|\det(x\mid y)|$ for the original one.
\end{proof}

\subsection{The Gram determinant}
\label{sec:gram_determinant}

\begin{lemma}[Gram determinant of a square matrix]
\label{lem:gram_determinant_square}
For $M \in \mathrm{Mat}_{n\times n}(\mathbb{R})$,
\[|\det(M)| = \sqrt{\det(M)\det(M\transp)} = \sqrt{\det(M\transp)\det(M)}
  = \sqrt{\det(M\transp M)}.\]
\end{lemma}

\begin{proof}
The first equality is $|\det M| = \sqrt{(\det M)^2}$ together with $\det(M\transp)=\det(M)$;
the last is multiplicativity of the determinant, $\det(M\transp)\det(M)=\det(M\transp M)$.
\end{proof}

% Correction: Claude Opus 5 (Ultracode) --- read "the volume ... is given by the Gram
% determinant", which contradicts the ainote immediately below fixing this document's convention:
% the bare $\det(M\transp M)$ is the Gram determinant and its *square root* is the volume. The
% same conflation was repeated at three later sites in this chapter and is corrected there too.
The quantity $\det(M\transp M)$ is called the \newterm{Gram determinant}
\germanfor{Gram-Determinante} of $M$. So, by \cref{prop:area_parallelogram_r2}, the
volume of an $n$-parallelotope in $\mathbb{R}^n$ is given by the \emph{square root} of the Gram
determinant, in analogy to $\lvert\det M\rvert$ itself.

% Supplement: Analysis_II_Script_v1.pdf, p. 119
\begin{ainote}
Terminology diverges from the official script here, worth flagging since it is exactly the kind
of thing that costs a mark in an exam. The script's Definition 13.55 calls $\sqrt{\det(L\transp
L)}$, the square root and so the \emph{volume} itself, the \qt{Gram determinant}. This
document instead follows the more standard usage and calls the bare $\det(M\transp M)$ the Gram
determinant, with the square root then giving the volume (\cref{def:volume_gram_determinant}
below). Every downstream formula agrees regardless of which name is attached to which quantity,
so this is terminology only, not a mathematical disagreement.

The script's own statement of Definition 13.55 is, independently, internally inconsistent: it
introduces a linear map $L:\mathbb{R}^m\to\mathbb{R}^n$ with $m<n$, but then writes the formula
using an undefined $d$ (\qt{the $d\times d$ matrix $L\transp L$}) where the surrounding text
means $m$; and the prose describes \qt{the determinant of the \dots matrix $LL\transp$} while the
displayed formula correctly uses $L\transp L$. Only $L\transp L$ (size $m\times m$) is invertible
for $m<n$, since $LL\transp$ is $n\times n$ and singular. The displayed formula is therefore the
reliable half of the definition, and the surrounding prose is not.
\end{ainote}

\begin{ainote}
The point of restating $|\det M|$ this way is that the middle expression, $\det(M\transp
M)$, still makes sense once $M \in \mathrm{Mat}_{n\times m}(\mathbb{R})$ is rectangular:
$M\transp M \in \mathrm{Mat}_{m\times m}(\mathbb{R})$ is square even when $M$ is not, so its
determinant is always defined, whereas $\det(M)$ itself is not. This is what lets the
notion of \qt{volume spanned by vectors} survive dropping the requirement $m=n$. A set of
fewer than $n$ vectors in $\mathbb{R}^n$ still spans a lower-dimensional parallelotope, and
that parallelotope still has an $m$-dimensional volume, even though no square matrix and no
ordinary determinant is available to compute it with.
\end{ainote}

\begin{definition}[Volume of a parallelotope via the Gram determinant]
\label{def:volume_gram_determinant}
For $m$ vectors $v_1,\dots,v_m \in \mathbb{R}^n$, the \newterm{area}/\newterm{length}/
\newterm{volume} of the parallelotope spanned by them is
\[\sqrt{\det\big(\langle v_i,v_j\rangle\big)} = \sqrt{\det(V\transp V)},
  \qquad V := (v_1\mid v_2\mid\cdots\mid v_m) \in \mathrm{Mat}_{n\times m}(\mathbb{R}).\]
\end{definition}

\begin{example}[$m=1$: the length of a vector]
\label{ex:gram_m1_length}
For $m=1$, $\sqrt{\det(v_1\transp v_1)} = \sqrt{\langle v_1,v_1\rangle} = |v_1|$,
which is the length of $v_1$.
\end{example}

% Source: Corsin Nick/Class Notes/Week 9.pdf, p. 4
\begin{example}[$m=2,\ n=3$: area of a parallelogram in $\mathbb{R}^3$]
\label{ex:gram_m2_n3_area}
For $m=2$, $n=3$, we hope that the Gram determinant of
\[M = \begin{pmatrix} x_1 & y_1 \\ x_2 & y_2 \\ x_3 & y_3 \end{pmatrix}\]
gives the area of the parallelogram spanned by $x$ and $y$ in $\mathbb{R}^3$. If $\theta$ is
the angle between $x$ and $y$, this should be $A = |x|\,|y|\,|\sin\theta|$.

% Sonnet 5 (Medium) --- FIG-W09-02
% Correction: Gemini 3.7 Flash (High) --- dropped altitude was vertical rather than perpendicular
% to y (angle 81.6 deg); now drops perpendicularly to proj_y(x) = (2.79, 0.41) on the dashed line of y.
% Correction: Claude Opus 5 (Ultracode), found by rendering printed p. 314 --- that fix was right
% and incomplete. Three things follow it.
%   * The $|x|\sin\theta$ label had no plate, and the purple edge $XY \to Y$ runs straight through
%     it: that edge is $y = 0.25 + 0.6538(x-1.7)$, which crosses the label's own height $y = 1.05$
%     at $x = 2.92$, inside the label's span. Printed, the bar and the $x$ of $|x|$ came out
%     struck through. It takes the `lbl' plate from style.md. (The strikethrough predates the
%     altitude fix -- at the old position $(2.68,1.0)$ the same edge crossed at $x = 2.847$ --
%     but the label was moved without checking the page.)
%   * The $\theta$ arc annotates the very angle the altitude depends on and was left unaudited.
%     `(0.55,0.075) arc (5:33:0.6)' starts at radius $0.5551$, so TikZ centres the arc at
%     $(0.55,0.075) - 0.6(\cos 5^\circ, \sin 5^\circ) = (-0.048, 0.023)$, not at the origin; and
%     it spans $5^\circ$ to $33^\circ$ where the true directions are
%     $\arctan(0.25/1.7) = 8.37^\circ$ and $\arctan(1.7/2.6) = 33.17^\circ$. Restarted from
%     $0.6(\cos 8.37^\circ, \sin 8.37^\circ) = (0.5936, 0.0873)$, which puts the centre on $O$ and
%     the ends on $y$ and $x$. This is the idiom `02-metric-normed-inner-product/03-cauchy-schwarz'
%     already uses correctly (`arc (14:59:0.5)' against $14.04^\circ$ and $59.04^\circ$).
%   * A right-angle mark at the foot, so the perpendicularity the fix installed is visible rather
%     than merely true. Unit vector along $y$ is $(0.9893, 0.1455)$; the mark is drawn at side
%     $0.12$ from the foot back along $-y$ and up along the perpendicular $(-0.1455, 0.9893)$.
\begin{center}
\begin{tikzpicture}[>=Latex, scale=1.3,
  lbl/.style={fill=white, fill opacity=0.85, text opacity=1, inner sep=1.2pt,
              rounded corners=1pt}]
  \coordinate (O) at (0,0);
  \coordinate (X) at (2.6,1.7);
  \coordinate (Y) at (1.7,0.25);
  \coordinate (XY) at (4.3,1.95);
  \draw[orange, thick, ->] (O) -- (Y) node[below right, font=\scriptsize] {$y$};
  \draw[gray!55, dashed] (Y) -- (3.05,0.45);
  \draw[blue, thick, ->] (O) -- (X) node[above, font=\scriptsize] {$x$};
  \draw[purple, dotted, thick] (X) -- (XY) -- (Y);
  % Dashed, not dotted: the parallelogram's own edges above are thick dotted, and red against
  % purple is a distinction in hue alone, which style.md forbids relying on -- an e-reader renders
  % both as the same grey. The dash pattern, the right-angle mark and the label are the three
  % channels that now separate the altitude from the edge it runs beside.
  \draw[red, dashed, thick] (X) -- (2.79,0.41);
  \draw[red] (2.6713,0.3925) -- (2.6538,0.5112) -- (2.7725,0.5287);
  \node[lbl, text=red, font=\scriptsize, right] at (2.75,1.05) {$|x|\sin\theta$};
  \draw (0.5936,0.0873) arc (8.37:33.17:0.6);
  \node[font=\scriptsize] at (0.9,0.32) {$\theta$};
\end{tikzpicture}
\end{center}

And indeed:
\[
\begin{aligned}
\sqrt{\det(M\transp M)}
  &= \sqrt{\det\begin{pmatrix} \langle x,x\rangle & \langle x,y\rangle \\
    \langle x,y\rangle & \langle y,y\rangle\end{pmatrix}} \\
  &= \sqrt{\langle x,x\rangle\langle y,y\rangle - \langle x,y\rangle^2} \\
  &= \sqrt{|x|^2|y|^2 - |x|^2|y|^2\cos^2\theta} \\
  &= |x|\,|y|\,|\sin\theta| \qquad (=\ |x\times y|).
\end{aligned}
\]
\end{example}
